\subsection{Final fitting, calibration and monotone prediction}
After selection, a new GBT-B model was fitted on all 34,541,470 valid
2017--2022 base windows (207,248,820 threshold-expanded rows). An isotonic
calibrator was fitted on a deterministic 10\% sample of the 2023 calibration
partition:
671,768 of 6,719,274 base windows, expanded to 4,030,608 threshold examples.
On this calibration sample, ECE decreased from 0.02264 before calibration to
0.00165 after calibration. The final GBT fit took
12,286.9 seconds (about 3 hours 25 minutes).

After calibration, equal-weight pool-adjacent-violators algorithm (PAVA) was
applied to the six anchor probabilities
$[0.2,0.5,1,2,5,10]$~mm. The projection minimizes squared displacement subject
to a non-increasing probability sequence. Any requested threshold inside
$[0.2,10]$~mm is estimated by linear interpolation between projected anchors;
out-of-domain thresholds are rejected.

\begin{table}[htbp]
\centering\small
\caption{Held-out 2024 threshold monotonicity and weighted calibration error.
Violation rates use 31,872,380 adjacent threshold pairs.}
\label{tab:monotonicity}
\begin{tabular}{l r r r}
\toprule
Prediction layer & Weighted ECE & Violation rate (\%) & Maximum increase \\
\midrule
Raw selected GBT & 0.022701 & 0.019603 & 0.016010 \\
Isotonic calibrated & 0.001670 & 0.009739 & 0.035351 \\
Final PAVA output & 0.001670 & \textbf{0.000000} & \textbf{0.000000} \\
Climatology & 0.001425 & 0.000000 & 0.000000 \\
\bottomrule
\end{tabular}
\end{table}

Isotonic calibration alone did not enforce ordering across different rainfall
thresholds. PAVA did: the final anchor probabilities had zero strict
violations and zero maximum increase across all 31,872,380 adjacent pairs.
Linear interpolation also remained non-increasing across the five requested
intermediate thresholds, covering 25,497,904 additional adjacent threshold
pairs with zero violations (Table~\ref{tab:monotonicity}).

